%%This is a very basic article template.
%%There is just one section and two subsections.
\documentclass[12pt,  a4paper]{scrartcl}
\usepackage[latin1]{inputenc}
\usepackage[T1]{fontenc}
\usepackage{graphicx}
\usepackage{epstopdf}
\usepackage{avant}
\usepackage[ngerman]{babel}
\usepackage[a4paper, top=3cm, bottom=3cm, left=2.5cm,
right=1.5cm]{geometry}
\usepackage{fancyhdr}
\usepackage{listings}
\usepackage{paralist}
\usepackage{color}
\usepackage{wrapfig}
\usepackage{textcomp}


\definecolor{darkblue}{rgb}{0,0,.5}
\usepackage{hyperref}
\hypersetup{colorlinks=true, breaklinks=true, linkcolor=darkblue,
menucolor=darkblue, urlcolor=darkblue}
\usepackage{breakurl}

\usepackage{url}

\usepackage[cc]{titlepic}
\usepackage{graphicx}
\titlepic{\includegraphics{images/Xtext-logo-dark.eps}}


\usepackage{tabularx}

\addtokomafont{sectioning}{\color{darkblue}}
\addtokomafont{subsubsection}{\color{darkblue}}

\usepackage{lmodern}


% Setzen des Headers
\setlength{\headwidth}{\textwidth}
\fancyhead[L]{\raisebox{-0.05cm}{\includegraphics[height=0.6cm]{images/itemis_logo_4c.eps}}}
\fancyhead[R]{\fontfamily{pag}\selectfont{\textbf{LWC13 -
\raisebox{-0.09cm}{\includegraphics[height=0.5cm]{images/Xtext-logo-dark.eps}} Submission}}}

% Wir wollen eine "Fancy" Ausgabe
\pagestyle{fancy}

% Absatzformatierung
\frenchspacing
\parindent 0pt      % kein Einrcken bei neuem Absatz
\parskip 10pt    	% Abstand zwischen den Absaetzen

% Listingsformatierung
\lstloadlanguages{Java, XML}
\lstset{
    language=Java,
    showspaces=false,
    showstringspaces=false,
    basicstyle=\small\ttfamily\footnotesize,
    columns=fullflexible,% typewriter font look better with fullflex
    keywordstyle=\color[rgb]{0.627,0.165,0.467},
    stringstyle=\color[rgb]{0.286,0.172,0.980},
    commentstyle=\color[rgb]{0.294,0.565,0.478},
    tabsize=2,
    numbers=left,
    numberstyle=\scriptsize,
    breaklines=true,
    identifierstyle=\ttfamily,
    backgroundcolor=\color[rgb]{0.9, 0.9, 0.9},
    rulesepcolor=\color[rgb]{0.3, 0.3, 0.3},
    captionpos=b,
    frame=shadowbox,
    otherkeywords={@SuppressWarnings,@Singleton,@Check,@Inject}
}
\lstdefinelanguage {Xtext}[]{Java} {
    morekeywords={returns,generate,as,import,with,grammar}
}
\lstdefinelanguage {Xtend}[]{Java} {
    morekeywords={def,dispatch,override,val,var,as,IF,ENDIF,FOR,ENDFOR,BEFORE,AFTER,SEPARATOR,typeof,@Inject,extension},
    literate={«}{\guillemotleft}{1}
    		 {»}{\guillemotright}{1} 
%    otherkeywords={+=,=,\{,\}}
}
\lstdefinelanguage {Mwe2}[]{Java} {
    morekeywords={module,var,component,bean}
}
\lstdefinelanguage {dmodel}[]{Java} {
    morekeywords={datatype,entity,op,\:}
}
\lstdefinelanguage {QL}[]{Java} {
    morekeywords={form,\:}
}
\lstdefinelanguage {QLS}[]{Java} {
    morekeywords={page, section, uses, question, font\-style, font\-color,
    font\-family, font\-weight, widget, navigation, CheckBox, Radio,\:, \-},
    alsoletter={-}
}
\lstdefinelanguage {Xbase}[]{Java} {
    morekeywords={val,var,\:}
}
\lstdefinelanguage {instances}[]{Java} {
    morekeywords={\=}
}


%
%%% Trennungsregeln fuer das gesamte Dokument.
%
\hyphenation{language-work-benches domain-model
Generate-Domain-model-Language}

\renewcommand{\baselinestretch}{1.25}

\begin{document}
\begin{titlepage}

\begin{center}

% Upper part of the page
\raisebox{-6cm}{
  \includegraphics[width=10cm]{./images/Xtext-logo-dark.eps}
}
\\[3cm]


% Title
\rule{\linewidth}{0.7mm}
\\[0.4cm]


    \fontfamily{pag}\selectfont
    \LARGE{\color{darkblue}Language Workbench Competition 2011 \\
    Xtext Submission} \\[0.5cm]

    \color{black}
    \large{Version: 1.3 DRAFT - \today}
    \\[0.4cm]
\rule{\linewidth}{0.7mm}
\\[1.5cm]
\normalfont

% Author
\fontfamily{pag}\selectfont
Karsten Thoms (itemis)
\normalfont

\vfill

% Bottom of the page
%{\large \today}

\end{center}

\end{titlepage}


\newpage

\section*{Abstract}
The Language Workbench Competition 2011 (LWC11) is an initiative created by a
group of experts at the CodeGeneration 2010
conference\footnote{\url{http://www.codegeneration.net/cg2010/}}. The aim is to
set a common task\footnote{see \url{http://www.languageworkbenches.net/} for the
detailed description of the LWC11 competition and other submissions}
for Language Workbenches\footnote{\url{http://martinfowler.com/articles/languageWorkbench.html},
\url{http://blog.efftinge.de/2007/11/definition-of-term-language-workbench.html}}
which is implemented with the different existing alternatives in a comparable
way. This document describes in detail how the task is solved with
Xtext\footnote{\url{http://www.xtext.org}}. Xtext is one of the most well known
Language Workbenches and part of the Eclipse Modeling
Project\footnote{\url{http://www.eclipse.org/modeling}}.


\section*{Testimonial}
This project was developed with the help of some colleagues at itemis, and I am
grateful for their help:

\begin{compactenum}
  \item \emph{Sven Efftinge} contributed the Instance DSL. Without his knowledge
  it would be very hard to implement this project already with Xtext 2.0 in this
  pre-release state.
  \item \emph{Rainer Klute} has written the initial version of the introduction.
  Although at that time the task was realized with Xtext 1.0.2 and many places
  have been reworked, major parts of the Overview and Phase 0.1 descriptions are
  based on his contribution.
  \item \emph{Karsten Nolte and Alexander Hannweg} initially transformed the
  Wiki pages to LaTeX.
  \item \emph{The Xtext Development Team} was always responsive for questions
  , provided with the Domainmodel example a good template to create major parts
  of the implementation and was eager to fix detected bugs.
  \item Detailed feedback helped to improve the document. My thanks for
  reading the document carefully and providing input go to
  \footnote{Not all from itemis}
  : \emph{Markus V�lter, J�rg Reichert, Joel W. Denton}
\end{compactenum}

\newpage
\section*{Document History}

\textbf {Version 0.1 - 2013-01-28}
\begin{compactitem}
\item Initial creation \\
\end{compactitem}



\newpage

\renewcommand{\contentsname}{Table Of Contents}
\setcounter{tocdepth}{4}
\tableofcontents

\newpage

\normalsize

\section{Introduction}

\subsection{Task Description}
The {\href{http://www.languageworkbenches.net/images/5/53/Ql.pdf}{LWC13 task}}
is to implement a DSL for questionnaires (Questionnaire Language, QL), which
basically allows the definition of forms with questions.


\subsection{Technology Stack}
This tutorial expects that you are somehow familiar with Java and Eclipse and
have heard about \url{EMF} and how it works in general before. We start almost at the
beginning, but not quite :-) 

We will use Xtext 2.3.1, which is at the moment of writing the latest official
release.
Xtext 2.4 is in preparation and will be released with Eclipse Kepler in June
2013\footnote{\url{http://wiki.eclipse.org/Kepler/Simultaneous_Release_Plan}}.
The solution approach described here would work also with any version
of Xtext >= 2.0, but the API might differ slightly, so there is no guarantee
that each codeline printed here would work exactly with all versions. For better
reproduction it is highly recommended to use the versions mentioned above.

For Code Generation we will use the language Xtend, which itself is based on
Xtext. Xtend makes use of a common expression language shipped with Xtext called
Xbase. The languages developped here will also be based on Xbase, but more on
this later.

\subsection{Xtext Overview}

This overview will give you a rough idea about what
Xtext\footnote{\url{http://www.xtext.org}} is all about. We will then dive into
the details and work on a small project.

In a nutshell, Xtext is a workbench to create and work with textual
domain-specific languages (DSLs). It comes as a feature (set of plugins) to the
popular Eclipse IDE.

The first thing you will want to do is to create your own domain-specific
language (DSL) and specify a \emph{grammar} for it. The grammar file is a plain
text file with ``\texttt{.xtext}'' filename extension, and the grammar within is
defined with a BNF like syntax. While you can use any text editor to modify it,
Xtext gives you a specialized editor for grammar files. It is aware of the Xtext
language, gives you syntax coloring, code completion, and more. To get a first
impression see the screenshot of the Xtext grammar file, opened with the Xtext
grammar editor, below. It is not required to fully understand the content yet,
this will be discussed in the next chapter in detail.

\includegraphics[width=16cm]{images/overview-xtext-grammar.png}

An example for DSL is one that allows you to define entities like ``Person'',
``Car'', ``Book'', and so on.
An entity has properties, e.g. a
Person has a name, a gender, and a date of birth. A Book has a title, one or
more authors, and an ISBN number.

A textual DSL model could look like this, but you could also imagine other
syntaxes:
\begin{lstlisting}
entity Person {
  name : String
  gender : m
  birthday : Date
}

entity Book {
  title: String
  authors: Person[]
  isbnNumber: String
}
\end{lstlisting}

Note that the Property \texttt{authors} is of type \texttt{Person}, so there can
be references between entities. In the Xtext grammar file you specify how you
want to define entities and their properties.

Once you have completed your language, you can do that: define some entities,
say ``Book'' and ``Person'', together with their respective properties and with
proper references between them. The nice thing is that Xtext not only gives you
a syntax-driven editor for editing grammar files. Additionally it generates an
editor that is specific to the language you have defined. It knows about your
language's keywords and where to place them, it knows about all the syntactical
constructs you have made up in your grammar, it includes all the nice stuff like
syntax coloring, code completion, validation, and more. For example, if you are
at some point where a reference to another entity must be inserted, your DSL
editor shows you all the references that would be valid here -- according to
your language rules -- and lets you choose among them. All in all, using the DSL
editor generated by Xtext, it is quite easy to establish a text file that adhers
to your DSL.

Depending on your language's type, you could call this text file e.g. a model, a
document, a program, or whatever. We will refer to DSL files here as
\emph{models} (files).

Consider now that you have created a model. What can you do with it? A typical
requirement is to generate an implementation of it in a language like Java, C++,
or XML. Or a graphical representation. Or something quite different. This is
where code generation comes in. Xtext creates a skeleton
code generator for you. Typically you use that code generator as a starting
point to produce e.g. Java source code, documentation in, say, DocBook or Wiki
format, overview graphics using GraphViz, or any other stuff you need. Xtext
offers special support for textual output formats, but it is also possible to
generate binaries.

This was only a short outline of some prominent Xtext aspects. It is by far not
everything Xtext can do for you, but it should suffice for now. The next
chapters will show you in more detail how to work with Xtext.

\subsection{Installing Eclipse and Xtext}

Xtext is a SDK for the \href{http://www.eclipse.org/}{Eclipse} IDE. To
install it you have two options:

\begin{compactitem}
    \item You can download Xtext separately and install it in your Eclipse
    instance.
    \item You can download a specially-crafted complete Eclipse distribution
    which has Xtext prepackaged already.
\end{compactitem}

We will take the latter approach here and describe the individual steps:

\begin{compactenum}
    \item Go to the \href{http://xtext.itemis.com/xtext/language=en/36553/downloads/}{Xtext download
    page}. Here you can get Eclipse 4.2.x (Juno) including Xtext 2.3.1
    along with some tools Xtext depends on. The latter are subsumed here under
    ``Xtext'' for simplicity.
    If you want you can download also a distribution
    which is already bundled with Eclipse 4.3.0 Kepler, but be aware that this
    is not finalized until end of June 2013.
    \item The Eclipse/Xtext distribution is available for multiple platforms.

    \begin{compactenum}
      \item
      \href{http://download.itemis.com/distros/eclipse-SDK-4.2-Xtext-2.3.1-linux-gtk-x86_64.tar.gz}{Linux GTK x86 64 bit}
      \item
      \href{http://download.itemis.com/distros/eclipse-SDK-4.2-Xtext-2.3.1-linux-gtk.tar.gz}{Linux GTK x86 32 bit}
        \item
        \href{http://download.itemis.com/distros/eclipse-SDK-4.2-Xtext-2.3.1-macosx-cocoa-x86_64.tar.gz}{Mac OSX x86 64 bit}
        \item
        \href{http://download.itemis.com/distros/eclipse-SDK-4.2-Xtext-2.3.1-win32-x86_64.zip}{Windows 64 bit}
        \item
        \href{http://download.itemis.com/distros/eclipse-SDK-4.2-Xtext-2.3.1-win32.zip}{Windows 32 bit}
    \end{compactenum}

    \item Unpack the downloaded archive file in a directory of your choice.

        Example (Linux):

\begin{lstlisting}
  cd /opt/local
  gzip -dc /download/eclipse-SDK-4.2-Xtext-2.3.1-linux-gtk-x86_64.tar.gz | tar
  xvfp -
\end{lstlisting}

        The archive will be extracted to a new directory named \texttt{eclipse}. Before
        unpacking the archive, please ensure that there is no subdirectory named
        \texttt{eclipse} yet! Different operating systems may require different unpacking
        methods.\footnote{On Windows do not unpack it into a deep directory,
        since this might cause troubles with long path names.}

    \item Start Eclipse by running the \texttt{eclipse} executable in the newly-created
    \texttt{eclipse} directory.
\end{compactenum}

\subsection{Workspace Setup}

Before we begin, start Eclipse and set up a fresh workspace.

Some settings should be done. Open the workspace settings:

\begin{compactitem}
    \item Windows: Window / Preferences
    \item Mac: Eclipse / Preferences
\end{compactitem}

\paragraph{Workspace Encoding}
$\;$ \\
File encoding is important to some type of files. It is better that the
workspace is set to a common encoding to avoid any platform specific encoding.
By default the workspace is using platform encoding, which is Cp1252 on Windows
and MacRoman on Mac. We will use ISO-8859-1 as a common encoding here.

\begin{compactitem}
    \item Open Eclipse Preferences and go to \emph{General / Workspace}
    \item Change setting \emph{Text file encoding} to \emph{Other /
    ISO-8859-1}
\end{compactitem}

\paragraph{Launch Operation}
$\;$ \\

\begin{compactitem}
    \item Open Run/Debug / Launching
    \item Change ``Launch Operation'' to ``Always
    launch the previously launched application''
\end{compactitem}
This will allow you re-running the previous launched application by just
pressing the Run or Debug button in the Eclipse toolbar, or using keyboard
shortcuts. The default settings does not always do what you want.


\input{introduction_050_survey_application.tex}

\section{Developing the Questionnaire Language}

\section{Developing the Questionnaire Language}

Step 1: Set up initial language implementation
- Create Project
- Define Grammar
- Generate Implementation

\subsection{Generate Language Implementation}

Now that the initial grammar of the language has been defined it is time to test
the language. Xtext ships with a code generator which generates all the glue
code needed for the language implementation.

To generate the code, we need to execute the generator workflow
\texttt{GenerateQlDsl.mwe2}. For this, select the workflow file, open the context menu
and select \emph{Run As / MWE2 Workflow}.

The generator will print some information to the Console, and finally it should
print \texttt{``Done.''}.

\begin{lstlisting}
0    [main] INFO  lipse.emf.mwe.utils.StandaloneSetup  - Registering platform
uri ....
   ...
13727 [main] INFO  .emf.mwe2.runtime.workflow.Workflow  - Done.
\end{lstlisting}

After successful execution the projects will be filled with implementation code.
Code that will be regenerated each time the generator is executed will go to the
source folder \texttt{/src-gen} (in all three projects), whereas code generated to
\texttt{/src} will be generated only once as skeleton. It is safe to edit these
classes.

Xtext follows the Generation Gap
Pattern\footnote{\url{http://heikobehrens.net/2009/04/23/generation-gap-pattern/}}:
Generated code is based on the Xtext API. Manual code is separated from
generated code. Often manual classes are derived from generated classes to allow
overriding of generated code or adding functionality.

Investigate the generated code a bit. Some pieces to mention:
\paragraph{Runtime Project - folder \texttt{src}}
\begin{itemize}
\item The class \texttt{QlDslRuntimeModule} is a Guice configuration.
Guice\footnote{\url{http://code.google.com/p/google-guice/}} is a famous
Dependency Injection\footnote{\url{http://en.wikipedia.org/wiki/Dependency_injection}} framework in Java. 
Xtext makes heavy use of Dependency Injection, which in turn allows to exchange
nearly every bit of the framework for customizing or to work around limitations,
if necessary, without the need to change the framework itself.
\item Class \texttt{QlDslStandaloneSetup} is needed when using the language in
``standalone mode'', i.e. without an Eclipse environment. Eclipse plugins, like
Xtext and the language plugin, usually need an OSGi container as execution
environment. Xtext is designed to be executable without the need to be
deployed into an OSGi container, but for this certain registrations are
required which an OSGi container would usually provide automatically. This is
especially useful when Xtext based languages are used in build environments or
other IDEs.
\item Class \texttt{QlDslFormatter} allows the implementation of a declarative
code formatter for the DSL.
\item File \texttt{QlDslJvmModelInferrer.xtend} is a class implemented with the
Xtend language. The JVM Model Inferrer will play an important role later when we
introduce expressions and code generation.
\item Class QlDslJavaValidation allows the implementation of validation rules
for the DSL.
\end{itemize}


\paragraph{Runtime Project - folder \texttt{src-gen}}
\begin{itemize}
  \item The Ecore metamodel is generated to file \texttt{QlDsl.ecore}.

\includegraphics[width=16cm]{./images/chapter01/QlDsl_ecore.png}

  \item The Java implementation code for the metamodel can be found in the
  package \newline\texttt{org.eclipse.xtext.example.ql.qlDsl}.
  \item The package \texttt{org.eclipse.xtext.example.ql.parser.antlr.internal}
  contains an \newline ANTLR3\footnote{\url{http://www.antlr.org}} grammar and
  the Lexer and Parser classes generated from it.
\end{itemize}


\section{Testing the Questionnaire Language}

- Create Launch Config
- Start Runtime
- Create Test Project

\section{Xbase}

- Describe Xbase
- Examples for Expressions
- Hint: Xbase Tutorial Projects

Extending Xbase: http://koehnlein.blogspot.de/2011/07/extending-xbase.html

\subsection{Including Expressions into the QL Language}

- Show final example (explain what is missing?)
- Derive Grammar from Xbase
- introduce Money datatype => operation overloading
- Generate Implementation

\subsection{JVM Model Inference}

For languages using Xbase it is necessary to tell Xtext, how to map concepts of a language to a Java model. In our example,
a Form could be mapped to the Type concept, while Questions are the fields of a class. By doing this, elements of the language
can be made available in expressions. Further, it allows that model elements are linkable where Java types are expected, without
necessarily generate a Java class.

The derivation of the Java model for language concepts is the responsibility of the JVM Model Inferrer, which is a class that implements
the \href{http://download.eclipse.org/modeling/tmf/xtext/javadoc/2.3/org/eclipse/xtext/xbase/jvmmodel/IJvmModelInferrer.html}{IJvmModelInferrer} interface.
A skeleton has already been generated into package org.eclipse.xtext.example.ql.jvmmodel. The file QlDslJvmModelInferrer.xtend is a class
written with Xtend.


The mapping that has to be implemented for the Questionnaire DSL should be as follows:
\begin{enumerate}
  \item Each Form instance is mapped to a JvmDeclaredType (which is the common concept for Java classes and interfaces).
  The type's name is simply the form name, and the target package is forms.
  \item Each Question of a Form is mapped to a JvmField, which is added as member of the declared type
  \item For each Question accessor methods for the field are generated. The field gets only a Setter, when the value of the Question is
  not computed by an expression. If the field is computed, the content of the getter has to compute the result.
  \item For
\end{enumerate}

- Describe the role of JVM Types
- Describe the role of the JVM Model Inferrer
- Implement JVM Model Inferrer



\section{Developing the Code Generator} \label{chp:QLCodeGenerator}

\subsection{Reference Implementation}

Before implementing a code generator one has to know what the target code is. Therefore a reference implementation has been
developed which was coded to large degree manually. From this reference code the templates can be derived. Also this is a
manual step.

We use a Java Server Faces (JSF) based application, which can be deployed on any Java Web Application server like Tomcat, JBoss etc.

The following screenshot shows the structure of the web application project. The application is available for download from the project 
homepage (\href{http://lwc13-xtext.eclipselabs.org.codespot.com/files/JSF-QL-1.0.zip}{JSF-QL-1.0.zip})

\includegraphics[width=10cm]{./images/chapter02/referenceimpl_projecttree.png}

Large parts of the application are not derivable from the model, they build the skeleton of the project. This is:
\begin{itemize}
\item Custom types (\texttt{src/types/*})
\item Images (\texttt{WebContent/resources/default/img/*})
\item Page Templates (\texttt{WebContent/resources/default/templates/*})
\item Libraries (\texttt{WebContent/WEB-INF/lib/*})
\item Web Application Descriptor (\texttt{WebContent/WEB-INF/web.xml})
\end{itemize}

We will focus on the parts which are dependent on the QL model and thus subject of code generation. These artifacts are:
\begin{itemize}
\item Java Bean classes representing the state of a Form (\texttt{src/forms})
\item JSF enabled XHTML pages representing the presentation of a Form (\texttt{WebContent/forms/*})
\item Index page (\texttt{WebContent/index.xhtml})
\item Faces configuration (\texttt{WebContent/WEB-INF/faces-config.xml})
\end{itemize}



\subsection{Xtend}
 
Since we will use Xtend to write the code generator, this chapter describes some
basic concepts of the Xtend language. However, we will not cover all aspects of
Xtend in this section, for more information see also the official documentation\footnote{\url{http://www.eclipse.org/xtend/documentation.html}}.

Xtend is a statically-typed general purpose language similar to Java. Xtend uses
Xbase as core language which was already roughly explained in section \ref{sec:Xbase}.
Hence, the presented concepts of Xbase are also valid for Xtend. The main focus
of Xtend lies in providing a language that is more readable than Java in certain
situations. In the backgroud, Xtend compiles to Java code. Thus, it plays perfectly
together with Java, e.g. methods declared in Java classes can be called in Xtend
and vice versa. Several concepts of Xtend are especially beneficial when writing code
generators. 

\begin{lstlisting}[language=Xtend]
def someMethod() {
  var myQuestion = 'where do we go?'
  myQuestion = myQuestion.toFirstLower
  val myAnswer = 42
  'The answer for question '+myQuestion+' is: '+myAnswer
}
\end{lstlisting}

In Xtend, a method is defined with the \texttt{def} keyword. The return type is
optional and will be automatically inferred. Only when a method is recursively
invoking itself, a return type needs to be specified explicitly. The keyword 
\texttt{var} defines a variable; constant values are defined with \texttt{val}.
In Xtend - since it is based on Xbase - everything is an expression, meaning that 
it has a return type. The last expression of a method defines the return value
and also the return type of the method.

In line 3 a so-called extension method is used. Recall that the String class in 
Java does not provide a method \texttt{toFirstUpper}. Xtend allows for extending
closed types without changing (maybe you already guess where Xtend got its name from).
You can easily write your own extension for the String type for example:

\begin{lstlisting}[language=Xtend]
  def square(Integer input) {
  	input * input
  }
  
  def useExtension() {
  	16.square // returns 256
  }
\end{lstlisting}

What Xtend basically does is changing the syntax of how methods are called. Instead
of writing \texttt{toLastUpper(``Hello'')}, Xtend always offers the alternative to
use the first input parameter as receiver of the method call (\texttt{``Hello''.toLastUpper}).
This results in the syntax being more chained than nested which improves readability.
To use methods from another class as extension methods in your class, the field
defining an object of the other class needs to be marked with the \texttt{extension}
keyword. In our scenario we will use dependency injection like the following:

\begin{lstlisting}[language=Xtend]
class JSFGenerator implements IGenerator{
  @Inject extension IJvmModelAssociations
  ...
}
\end{lstlisting}

This statement simply allows to use the methods declared by the interface \texttt{IJvmModelAssociations}
in our generator class as extension methods on our objects. Which concrete implementation of the
interface is later actually called, is configured in the Guice module.

A further useful feature of Xtend is \emph{polymorphic dispatching}. Using the \texttt{dispatch}
keyword on multiple methods with identical signatures has the effect that the 
decision on which method should be called is based on the runtime type of the
target object. In contrast, Java binds methods at compile time based on the static
type of the target object. Since Xtend compiles to Java, polymorphic dispatching
is internally realized by a dispatcher method using a cascade of \textt{instanceof}
constructs. 

\begin{lstlisting}[language=Xtend]
  def dispatch doSomething(Integer input) {
  	// do something with integer
  }
  
  def dispatch doSomething(Float input) {
  	// do something with float
  }
  
  def useDispatching() {
    val Number x = 123
    val Number y = 0.1
    x.doSomething + y.doSomething
  }
\end{lstlisting}

// TODO gescheites beispiel überlegen und zwei sätze zu schreiben.

Xtend offers the possibility to define \emph{rich Strings} (also called \emph{templates})
which is especially useful when writing code generators. Rich Strings allow for
writing complex Strings with line breaks and indentations without the need for
concatenating special characters like \texttt{'\textbackslash n'} or \texttt{'\textbackslash t'}. A rich String
construct starts and ends with triple single quotes (\texttt{'''}). Within such
a String code pieces which themself return a String can be inserted (surrounded by guillemots \texttt{«»}).
There are also logical structures like \texttt{for} loops or \texttt{if-else}
statements supported:

\begin{lstlisting}[language=Xtend]
  def htmlContent(List<String> contents) '''
  	<html>
  	  <body>
  	    «FOR content: contents BEFORE '<p>' SEPARATOR '<br/>' AFTER '</p>'»
  	      «IF content.length > 10»
  	        Large Content: «content.toFirstUpper»
  	      «ELSE»
  	        Small Content: «content.toFirstUpper»
  	      «ENDIF»
  	    «ENDFOR»
  	  </body>
  	</html>'''
\end{lstlisting}

Note the special keywords in the \texttt{FOR} loop declaration: \texttt{BEFORE} 
and \texttt{AFTER} will be called once before and after the iteration, but only
if the loop will be iterated at least once. With the \texttt{SEPARATOR} keyword
a string which will be inserted between two iteratins can be specified. Calling 
the example method with ['hello', 'Some more text'] results in:

\begin{lstlisting}[language=Xtend]
<html>
  <body>
    <p>
    Small Content: Hello<br/>
    Large Content: Some more text
    </p>
  </body>
</html>
\end{lstlisting}

Last but not least, Xtend offers a more sophisticated \texttt{switch-case} statement
than Java does. The \texttt{break} statement as known from Java is implicit in Xtend.
Furthermore, switching based on Strings and even types is possible:

\begin{lstlisting}[language=Xtend]
  def getTypeName(Number input) {
  	switch (input) {
  		Integer: "It is an Integer!"
  		Float:   "It is a Float!"
  		default: "It is some other number type."
  	}
  }
\end{lstlisting}

Now as you know the basic concepts of Xtend, let's finally start writing the code
generator.

\subsection{Code Generator}

In this section you will learn how to implement the code generator for 
the target application. For simplicity the code generator templates are placed in the
org.eclipse.xtext.example.ql project in a sub-package generator. 
Usually it would be better to create a separate project which contains the generator,
since the language is independent from a single target platform. It would
be possible to create different code generators for different target platforms,
and it would be better to implement each of them as separate projects.

Generator templates in Xtend are implementations of the org.eclipse.xtext.generator.IGenerator
interface:

\begin{lstlisting}[language=Java]
public interface IGenerator {
	/**
	 * @param input - the input for which to generate resources
	 * @param fsa - file system access to be used to generate files
	 */
	public void doGenerate(Resource input, IFileSystemAccess fsa);
}
\end{lstlisting}

\subsubsection {Dispatcher template}

The code generator is invoked with a Resource instance, which holds a Questionnaire
instance. We have to generate multiple artifacts for each resource, so it is a common 
pattern to create a template class serves as entry point and dispatches to other
template classes to create the artifacts. Usually one template per artifact is
created.

Create the class Root.java in package org.eclipse.xtext.example.ql.generator:

\begin{lstlisting}[language=Java]
package org.eclipse.xtext.example.ql.generator;

import javax.inject.Inject;

import org.eclipse.emf.ecore.resource.Resource;
import org.eclipse.xtext.generator.IFileSystemAccess;
import org.eclipse.xtext.generator.IGenerator;
import org.eclipse.xtext.xbase.compiler.JvmModelGenerator;

@SuppressWarnings("restriction")
public class Root implements IGenerator {
  @Inject
  JvmModelGenerator jvmModelGenerator;

  public void doGenerate(Resource input, IFileSystemAccess fsa) {
    // dispatch to other generators
    jvmModelGenerator.doGenerate(input, fsa);
  }
}
\end{lstlisting}

As a first generator to which a dispatch, we inject an instance of
JvmModelGenerator. This is a standard generator shipped with Xtext which
translates types inferred by the Jvm Model Inferrer to Java classes.
In our case, the Java Bean class for Forms are generated by the
JvmModelGenerator. 
 
Next, Xtext has to know that Root.java is the template that has to be invoked
as generator implementation. This has to be configured - you guessed right! - 
with Guice again. We need to add a configuration that binds the IGenerator
interface to the Root class.

Open class QlDslRuntimeModule and add this method:

\begin{lstlisting}[language=Java]
  @Override
  public Class<? extends IGenerator> bindIGenerator() {
    return Root.class;
  }
\end{lstlisting}

Now we are ready to add templates additional templates and register them in the
Root class.

\subsection{Testing the Questionnaire Application}
\label{sec:TestingQLApplication}

It's time to test the application we have developed so far. If your runtime
environment is still open you need to restart it. Just switch to the runtime
instance and press \emph{File / Restart}. See also section \ref{sec:TestingQL}
for more information on how to start your runtime evironment for testing. Once
the runtime environement is running you can test the code generator in the
QLText project created in section \ref{sec:TestingQL}. Code generation process
is triggered automatically on the fly when the QL model gets modified. The
generated artifacts are located in folder \texttt{WebContent/generated/forms}.
Into this folder also an \texttt{index.xhtml} file is placed which you can use
as starting point for the questionnaire application. Right-click on the index
file and press \emph{Run As / Run On Server}. In the next dialog choose the
Tomcat server which you should have configured in section
\ref{sec:IDEConfiguration} (if not, then now is the time to do it) and press
\emph{Next}. If your test project is not already in the \emph{Configured} area
then it needs to be moved there by using the \emph{Add >} button. Finally press
\emph{Finish} to start the application in the integrated browser within eclipse.
You can also copy the link and paste it in a web browser of your choice. The
result should be similar to the one on the next screenshot:

\begin{center}
\includegraphics[width=15cm]{./images/chapter02/questionnaireApplication.png}
\end{center}
 
Now it's show time! You can modify your QL model and save it. The XHTML files
will be regenerated on the fly. The server needs some seconds to reload the
backing beans. Afterwards you can press the \emph{Refresh} button and see the
result immediately.



\section{Layout and Styling Language (QLS)}

\subsection{The Language QLS}

In this chapter we will describe how the optional task to define a language 
for styling and layout information can be accomplished in our language workbench.
The language QLS should allow for defining the following information:
\begin{itemize}
  \item Grouping of question forms into pages, sections and subsections
  \item Navigation links between pages
  \item Styling of question texts by defining font style, font type and color
  \item Defining the appearance of a question by specifying the widget type to render
\end{itemize}

The following shows an example model that we want to be able to define with QLS:

\begin{lstlisting}[language=QLS]
page HouseOwningPage {
	section house uses Box1HouseOwning {
		question hasBoughtHouse [widget: CheckBox]
		question valueResidue [
			font-style:  "italic"
			font-wight:  "bold" 
			font-color:  "#2233FF"
			font-family: "Arial"
		]
	}
	section garage uses GarageOwning {
		question hasBoughtGarage [widget: Radio["Yepp", "Nope"]]		
	}
	navigation {
		CarOwningPage
	}
}

page CarOwningPage uses CarOwning {
	// ...
}
\end{lstlisting}

This example model defines two pages. The first page consists of two sections,
one for the question form \emph{Box1HouseOwning} which was defined in chapter 
\ref{chp:QL} and one for a further form called \emph{GarageOwning}. The form to
be rendered inside a page or a section is specified by the \texttt{uses} keyword.
This definition must be unambiguously, e.g. if there is a form included by a page,
it is not allowed to refer to another form in one of its containing sections.
This restriction is ensured by implementing a corresponding validation which will be covered
in chapter (TODO). In a section (or page) the styling information for the questions of the included
form can be defined as in lines 3-9 and 13. The \texttt{navigation} keyword allows to define the order in which the pages are to
be displayed by specifying which page should appear next.

Before we can write the corresponding Xtext grammar, we need to create the DSL 
projects for QLS as we have done for the Questionnaire language. For this, refer
again to chapter \ref{chp:CreateDslProjects} and replace all occurences of 
\texttt{ql} with \texttt{qls}. The project \texttt{org.eclipse.xtext.example.qls}
should then contain the grammar file \texttt{QlsDsl.xtext}. We define the content
of this file as the following:

\begin{lstlisting}[language=Xtext]
grammar org.eclipse.xtext.example.qls.QlsDsl with org.eclipse.xtext.common.Terminals

generate qlsDsl "http://www.eclipse.org/xtext/example/qls/QlsDsl"
import "http://www.eclipse.org/xtext/example/ql/QlDsl" as ql

QuestionnaireStyleModel:
	pages+=Page*;
	
Page:
	"page" name=ID ("uses" form=[ql::Form|ID])? "{"
		element+=PageElement*
		navigation=Navigation?
	"}"
;

PageElement:
	QuestionStyling | Section
;

QuestionStyling:
	"question" question=[ql::Question] styling+=StyleInformation?	
;

StyleInformation: {StyleInformation}
	"["(
		("font-style:" fontStyle=STRING)? &
		("font-weight:" fontWeight=STRING)? &
		("font-color:" fontColor=STRING)? &
		("font-family:" fontFamily=STRING)? &
		("widget:" widget=Widget)?
	)"]"
;

Widget: {Widget}
	widgetType=("Radio"|"DropDown"|"CheckBox"|"Text"|"Slider") ("[" labels+=STRING ("," labels+=STRING)* "]")?
;
  
Section:
	"section" name=ID ("uses" form=[ql::Form|ID])? "{"
		element+=PageElement*
	"}"
;

Navigation: {Navigation}
	"navigation" "{" (nextPage+=[Page|ID])+ "}"
;
\end{lstlisting}

After reading chapter \ref{sec:DefiningGrammar} you should be familiar with the
concepts of Xtext grammar definitions and understand most parts of the QLS grammar.
One new concept represented in the QLS grammar is the one of unordered lists:

\begin{lstlisting}[language=Xtext]
StyleInformation: {StyleInformation}
	"["(
		("font-style:" fontStyle=STRING)? &
		("font-weight:" fontWeight=STRING)? &
		("font-color:" fontColor=STRING)? &
		("font-family:" fontFamily=STRING)? &
		("widget:" widget=Widget)?
	)"]"
;
\end{lstlisting}

The styling information can be defined in an arbitrary order in which each kind 
of element (font style, font color and so on) may only occur at 
most once. The arbitrariness of the order is expressed by the \texttt{'\&'} between
the style elements. The optionality of each style information is again declared by
the question mark \texttt{'?'}.

A further noteworthy aspect is the handling of references. In Xtext, references 
to language elements are expressed by using squared brackets. An example for this are the  
references to existing pages in the navigation section:

\begin{lstlisting}[language=Xtext]
Navigation: {Navigation}
	"navigation" "{" (nextPage+=[Page|ID])+ "}"
;
\end{lstlisting}

Here, \texttt{nextPage} is a reference to an existing page description which may even
be defined in a different file. \texttt{ID} defines which attribute should be used
for the reference's name. Xtext automatically creates hyperlinks to the referenced
elements which can be enabled by holding \texttt{Ctrl} and moving the curser over the
reference or by pressing \texttt{F3}.

To define references to elements defined in other langauges than the own,
the reference needs to be qualified. An example in QLS is the reference to questions
or forms defined in a QL model. To express this on the grammar level, the QL DSL
needs first to be imported:

\begin{lstlisting}[language=Xtext]
import "http://www.eclipse.org/xtext/example/ql/QlDsl" as ql
\end{lstlisting}

Since the QL DSL is imported under the name \texttt{ql}, references can now be 
qualified by using a double colon (\texttt{::}):

\begin{lstlisting}[language=Xtext]
QuestionStyling:
	"question" question=[ql::Question] styling+=StyleInformation?	
;
\end{lstlisting}

This grammar rule allows for refering all question elements defined in
any QL model in our test project. However, since there is always an unambiguous
form specified for a page or a section, only the questions defined within this 
form should be referable. This is a typical scoping issue which needs to be 
handled in the scope provider of the QLS language. We already learned about scoping
in section \ref{sec:scoping}. The scope provider for the QLS language looks like
the following:

\begin{lstlisting}[language=Java]
public class QlsDslScopeProvider extends AbstractDeclarativeScopeProvider {

  public IScope scope_QuestionStyling_question(EObject context, EReference reference) {
    Form usedForm = getUsedForm(context);
    List<Question> allQuestions = EcoreUtil2.getAllContentsOfType(usedForm,
        Question.class);
    return Scopes.scopeFor(allQuestions);
  }

  private Form getUsedForm(EObject context) {
    if (context instanceof Section) {
      Section section = (Section) context;
      if (section.getForm() != null) {
        return section.getForm();
      }
    } else if (context instanceof Page) {
      Page page = (Page) context;
      if (page.getForm() != null) {
        return page.getForm();
      }
    }
    if (context != null) {
      return getUsedForm(context.eContainer());
    }
    return null;
  }
}
\end{lstlisting}

The method \texttt{scope\_QuestionStyling\_question} is always invoked whenever the visible
elements for the attribute \texttt{question} in the grammar rule \texttt{QuestionStyling} 
need to be computed. Here, we just use a recursive algorithm to get the form declaration
of the parent section or page which is mandatory. With the helper method \texttt{getAllContentsOfType}
all questions defined in this form are collected and a scope containing these questions is returned.

Now it's time to play around with the QLS language and to test its features. For this,
switch to the test project in the runtime environment (see also section \ref{sec:TestingQL})
and create a file with the file extension \texttt{.qls}.
\subsection{Adapt the Code Generator}

- describe which parts of the code generator need to be adapted to use the QLS features




\section{Additional Concepts}

\subsection{Validation}

As an optional task the LWC13 task requires the implementation of
analysis rules. Xtext provides a validation framework which integrates into the
EMF Validation
framework\footnote{\url{http://www.eclipse.org/modeling/emf/?project=validation}}.
The Xtext User Manual contains a Validation chapter that is
worth reading additionally\footnote{see
\url{http://www.eclipse.org/Xtext/documentation.html\#validation}}.

We will show in this section how the constraints defined in the task description
can be realized with Xtext.

\subsubsection{Extending the Java Validator class}
With the first translation of the Xtext grammar, the generator has already
created the necessary infrastructure to implement custom validation rules.
Look into the package \texttt{org.eclipse.xtext.example.ql.validation}, you will
find a Java class \texttt{QlDslJavaValidator}. The class extends the
\texttt{AbstractQlDslJavaValidator} class, which is regenerated each time the
grammar is translated. Thus, the generation gap pattern is applied here again.
It is safe to extend the \texttt{QlDslJavaValidator} class manually.

But instead of implementing the constraints in Java, we will use Xtend again.
For easier integration, our Xtend based validator class will be inserted into
the class hierarchy of QlDslJavaValidator.

Create an Xtend class \texttt{QlDslXtendValidator}, and extend it from
\texttt{AbstractQlDslJavaValidator}.

\begin{lstlisting}[language=Xtend]
package org.eclipse.xtext.example.ql.validation

import javax.inject.Inject
import org.eclipse.xtext.validation.Check
import org.eclipse.xtext.xbase.XFeatureCall
import org.eclipse.xtext.xbase.XbasePackage
import org.eclipse.xtext.xbase.jvmmodel.IJvmModelAssociations

import static extension org.eclipse.xtext.nodemodel.util.NodeModelUtils.*
class QlDslXtendValidator extends AbstractQlDslJavaValidator {
  @Inject extension IJvmModelAssociations
}
\end{lstlisting}

Derive QlDslJavaValidator from the the Xtend validator:

\begin{lstlisting}[language=Java]
public class QlDslJavaValidator extends QlDslXtendValidator {
  // do nothing here, rules are implemented in Xtend
}
\end{lstlisting}

\subsubsection{Constraint: Ensure order of questions}

The first constraint in the LWC13 task is defined so:

\begin{quote}
\emph{
Test for cyclic dependencies. For instance, the following snippet should be rejected:
}
\begin{lstlisting}[language=QL]
if (x) { y: "Y?" boolean }
if (y) { x: "X?" boolean }
\end{lstlisting}

\emph{The reason is that y will only be asked for when x is true, but x will only get a value when y is true. 
Of course such cyclic dependencies could occur transitively and nested in expressions. 
Another way of stating this check is: the ordering of questions should be
consistent with how the question variables are used in conditions and computed values. 
}
\end{quote}

The task allows two approaches to achieve the goal. We will choose the second
approach: Check that elements referred in expressions have been declared before
their usage.

Xtext does not enforce that elements are declared before they are used
somewhere. The referred names simply must be in the scope of the context, which
the current scope implementation already accomplishes. We could implement this
constraint also in the scope provider of the language by restricting the scope
to elements that have been declared before. However, we will implement this as a
semantic constraint in the validator class.

To implement the constraint we need to know the location in the model where a
Question element is declared and where it is called in an expression. Besides
the Abstract Syntax Tree Xtext maintains a second model, which represents the
document. This is the so-called \emph{Node Model}, and the class
\texttt{NodeModelUtils} provides some utility functions to navigate from an AST
element to the node model. Since Questions that are referred in expressions are
local to the Form, we can simply compare the offset in the document of the
declared Question and the feature call in an expression.

Xtext validation rules are implemented as methods which are annotated with
\texttt{@Check}. The method name does not matter. Check methods are expected to
have exactly one parameter, which is of the type of element that has to be
checked. The base class provides methods to create error messages. For the case
of this constraint, the necessary context object type is \texttt{XFeatureCall}. 

Now open the \texttt{QlDslXtendValidator} class again and add this
method\footnote{\url{https://gist.github.com/kthoms/5240455}}:

\begin{lstlisting}[language=Xtend]
  @Check
  def void check_featureDeclaredBeforeCall (XFeatureCall featureCall) {
    val featureSource = featureCall.feature.sourceElements.head
    val nodeFeature = if (featureSource != null) featureSource.node else featureCall.feature.node
    val nodeCall = featureCall.node
    if (nodeFeature != null) {
      if (nodeFeature.offset > nodeCall.offset) {
        error(featureCall.feature.simpleName+" must be declared before.",featureCall, 
        XbasePackage::eINSTANCE.XAbstractFeatureCall_Feature, "ERR_FEATURE_CALL_BEFORE_DECLARATION", null)
      }
    }
  }
\end{lstlisting}

After restarting the workbench the situation will be recognized as an error:

\includegraphics[width=6cm]{images/chapter04/validation_1.png}

\subsubsection{Constraint: Type conformance check}

Next, the LWC13 task requires checking the type conformance in expressions:

\begin{quote}
\emph{Type check conditions and variables: the expressions in conditions should be
type correct and should ultimately be booleans. The assigned variables should be
assigned consistently: each assignment should use the same type.
}\end{quote}

Here we have to do nothing, since this constraint is already implemented in
Xbase. This works thanks to Xbase's type inference mechanism. The Xtend language
makes heavy use of this nice feature, which makes it almost unneccessary to
declare types anywhere. For dynamic languages this is natural, since the actual
type is known and evaluated at runtime. Static typed language often lack this
feature.

\includegraphics[width=6cm]{images/chapter04/validation_2.png}

\includegraphics[width=6cm]{images/chapter04/validation_3.png}


\subsubsection{Testing validation rules}

It is easy to test validation rules in a runtime environment. However, we will
show how these rules can also be unit tested. Also herefore the Xtext framework
already contains the necessary infrastructure. Remember that Xtext has already
created a test plugin with the initial generator run? Now it is time to make use
of it.

Again, it is easier to create the test with Xtend. Xtend allows us to create
simple models inline with Rich Strings, pass the result to a parser, and
validate the result. With Xtend this is a one-liner.

\begin{lstlisting}[language=Xtend]
package org.eclipse.xtext.example.ql.validation.test

import javax.inject.Inject
import org.eclipse.xtext.example.ql.QlDslInjectorProvider
import org.eclipse.xtext.example.ql.qlDsl.Questionnaire
import org.eclipse.xtext.junit4.InjectWith
import org.eclipse.xtext.junit4.XtextRunner
import org.eclipse.xtext.junit4.util.ParseHelper
import org.eclipse.xtext.junit4.validation.ValidationTestHelper
import org.eclipse.xtext.xbase.XbasePackage
import org.junit.Before
import org.junit.Test
import org.junit.runner.RunWith

@RunWith(typeof(XtextRunner))
@InjectWith(typeof(QlDslInjectorProvider))
class QlDslValidationTest {
  @Inject extension ParseHelper<Questionnaire> parseHelper
  @Inject extension ValidationTestHelper
  
  @Before
  def void setUp () {
    parseHelper.fileExtension="ql"
  }
  
  @Test
  def void testValidation_CallBeforeDeclaration_expectError () {
    '''
    form Foo {
      if (x) { y: "Y?" boolean }
      if (y) { x: "X?" boolean }
    }
    '''.parse.assertError(XbasePackage::eINSTANCE.XFeatureCall, "ERR_FEATURE_CALL_BEFORE_DECLARATION", "must be declared before")
  }

  @Test
  def void testValidation_CallBeforeDeclaration_expectSuccess () {
    '''
    form Foo {
      x: "foo" boolean
      if (x) { a: "X?" boolean }
    }
    '''.parse.assertNoErrors
  }
  
  
  // Type check conditions and variables: the expressions in conditions should be type correct and should ultimately be booleans. 
  // The assigned variables should be assigned consistently: each assignment should use the same type.  
  @Test
  def void testValidation_ConditionTypeCheck_expectError () {
    '''
    form Foo {
      if ("foo".length) { a: "X?" boolean }
    }
    '''.parse.assertError(XbasePackage::eINSTANCE.XMemberFeatureCall, 
       "org.eclipse.xtext.xbase.validation.IssueCodes.incompatible_types", "Type mismatch")
  }
  
  @Test
  def void testValidation_ConditionTypeCheck_expectSuccess () {
    '''
    form Foo {
      if ("foo".length>1) { a: "X?" boolean }
    }
    '''.parse.assertNoErrors
  }

  @Test
  def void testValidation_AssignmentTypeCheck_expectFailure () {
    '''
    form Foo {
      a: "X?" boolean ("foo".length)
    }
    '''.parse.assertError(XbasePackage::eINSTANCE.XMemberFeatureCall, 
       "org.eclipse.xtext.xbase.validation.IssueCodes.incompatible_types", "Type mismatch")
  }

  @Test
  def void testValidation_AssignmentTypeCheck_expectSuccess () {
    '''
    form Foo {
      a: "X?" boolean ("foo".length>1)
    }
    '''.parse.assertNoErrors
  }
}
\end{lstlisting}

Basically, the class is a plain JUnit 4 class. We have to use a special
JUnit execution class, \texttt{XtextRunner}, and provide a language specific initializer
class, \texttt{QlDslInjectorProvider}.

Next, the class adds two extension classes. ParseHelper provides a parse()
method for char sequences. This will parse and validate the model. Afterwards
the observed issues can be asserted with the methods from the 

\begin{lstlisting}[language=Xtend]
@RunWith(typeof(XtextRunner))
@InjectWith(typeof(QlDslInjectorProvider))
class QlDslValidationTest { 
  @Inject extension ParseHelper<Questionnaire> parseHelper
  @Inject extension ValidationTestHelper
  ...
}
\end{lstlisting}

Now the test methods can be implemented. They are super simple:

\begin{lstlisting}[language=Xtend]
  @Test
  def void testValidation_CallBeforeDeclaration_expectSuccess () {
    '''
    form Foo {
      x: "foo" boolean
      if (x) { a: "X?" boolean }
    }
    '''.parse.assertNoErrors
  }
\end{lstlisting}
  
Due to the tight Java integration, the unit tests of the Xtend class can be
executed by running them through the context menu (\emph{Run As / JUnit Test}).

\includegraphics[width=12cm]{images/chapter04/validation_4.png}



\section{Closing Words}

Thank you for reading this document. We have been writing it with the intention
that it should provide easy access for first-time users of Xtext to this
powerful Language Workbench. This explains also the size of the document. We
could have written it shorter, if we would assume more background knowledge of
the potential readers, or if we would leave many things with just short
comments.

Xtext has grown over years, gaining more and more
experience from real-life projects that want to leverage the power of DSLs in
integrated environments. We could only touch the surface of Xtext and Xtend
here, and tried to choose the most simple solution. The QL assignment did fit
quite well into what Xtext can provide mostly out-of-the-box, but its real power
is unvealed when DSLs have non-standard requirements. Almost every peace in
Xtext is highly customizable and it is one of the most flexible frameworks we
know of. This flexibility comes to the price of complexity. Xtext is well
documented by the reference manual, several blogs show advanced concepts in
detail, and hundreds of projects solve different real-world requirements. Many
open projects exists which are worth to study. Learning Xtext does not mean
following a single tutorial, it needs training. The sources and the debugger are
valuable friends when solving issues. The community around Xtext is very helpful
and likely the largest of all language workbenches.

We hope that this document was helpful for you to understand some concepts of
Xtext and gave you the right level of abstraction to learn how to use this tool.
If we helped you to decide to use Xtext or Xtend in your project, then it was
worth the effort and we would be pleased to hear from you! 

\par
\par
\par

\emph{\large{Karsten Thoms, Johannes Dicks, Thomas Kutz}}

\emph{April 2013}



\end{document}